\section{Execution cost and existing measurements}
\label{sec:cost}
Ibex already has a substantial measurement corpus. The scale suite and its
published benchmark page provide observations across operators, row counts,
and thread configurations~\cite{ibexBenchmarks}. This section uses a small,
explicitly identified snapshot to ground the cost explanation. It separates
what those measurements show from the mechanisms that can explain a cost;
timing a query alone does not identify its bottleneck.

\subsection{What determines the amount of work?}
Row count is only one dimension. Let $n$ denote input rows, $s$ the fraction
surviving a filter, $w$ the bytes of demanded payload per row, $g$ the number
of groups, and $m$ the number of join-result rows. These are explanatory
quantities, not parameters of a claimed Ibex optimizer cost model.

\paragraph{Reading and selecting.}
Column demand limits which payloads must be read. A selective scan can reduce
payload materialization toward $snw$ bytes, but predicate evaluation still
has its own input cost. Source encodings and page boundaries determine whether
skipped rows also avoid decoding work. A small result therefore does not imply
a proportionally cheap scan. Cached columns change the work again: gathering
from existing buffers avoids source decoding but still reads and writes values.

\paragraph{Computing and moving values.}
A row-local expression such as doubling a price visits $n$ values and produces
$n$ outputs. A projection can share payloads, whereas a filter generally gathers
survivors into output buffers. Fusion can remove intermediate passes and avoid
gathering unused columns. Fixed-width output volume depends on rows and width;
string output also depends on the number of character bytes. Consequently,
operations with similar arithmetic can have different allocation and copying
costs. These mechanisms motivate a bandwidth hypothesis, but the published
wall times alone do not establish memory-bandwidth saturation.

\paragraph{Maintaining state.}
Aggregation combines row processing with state indexed by $g$ groups. More
groups, additional key components, and different key representations alter
lookup and state-management work. A hash join similarly pays for its build
index, probe lookups, and materializing $m$ matches. Duplicate keys can make
$m$ much larger than either input. Reducing output width helps the gather,
but does not eliminate the cost of identifying the matches.

\paragraph{Ordering and reducing output.}
Sorting a complete table and selecting its top 100 rows request different
amounts of output and can use different algorithms. Ibex's bounded-heap top-k
path performs selection before gathering its small result, rather than
materializing a full sorted table. Useful input ordering can also avoid
sorting or admit the sorted aggregate described in Section~\ref{sec:operations}.
Maintaining time order can thus have an execution benefit as well as a semantic
purpose, although time order and group-contiguous order are different facts.

\evidence{\src{src/runtime/lazy_table.cpp}, selected decoding;
\src{include/ibex/runtime/pipeline.hpp}, fused map descriptions;
\src{src/runtime/join_chunked.cpp} and
\src{src/runtime/aggregate_chunked.cpp}, state and output construction;
\src{src/runtime/runtime_entry.cpp}, the TopK execution-path description.
These are implementation-based explanations, not measured attribution of the
timings below to individual phases.}

\subsection{The measurement boundary}
The inspected \src{docs/benchmarks.html} snapshot was generated on
3 September 2026 at 13:08 UTC and records commit
\code{0d991d3677f2927eb40c77ba9f489ad6877baa24}. Its source is
\src{benchmarking/results/scales_aws_20260903T100845.csv}.
The published methodology describes AWS r7i.2xlarge instances with 8 vCPUs
(4 physical cores with two hardware threads each), 64 GB of memory, and one
instance per engine. The timings below are reported means from that run,
not fresh measurements of the current checkout.

For that archived snapshot, the Ibex wrapper selects resident input tables
and parses and lowers the query before the timed iterations. It warms up, then times
\code{runtime::interpret}, including construction of the materialized result.
This measures execution rather than file loading or C++ compilation. Separate
parse-inclusive and scan-inclusive harness modes have different boundaries
and must not be mixed into that interpretation. Physical execution decisions
made inside \code{interpret} remain inside the timer.
The revised in-memory harness includes parsing and lowering; its new results
must be distinguished from this historical snapshot.

\evidence{\src{docs/benchmarks.html}, embedded \code{PAYLOAD} metadata and
timings; \src{docs/methodology.html};
\src{benchmarking/gen_website.py}, \code{avg_ms} payload construction;
\src{tools/ibex_bench.cpp}, \code{run_benchmark}, especially the
\code{!include_parse} branch. The selected page cells were checked against
their source CSV. The page snapshot does not expose individual iteration
samples or the invocation's repetition count; harness defaults alone do not
establish that count for this archived run.}

\subsection{Representative observations}
All times in Table~\ref{tab:cost-snapshot} are milliseconds. ``Default'' is
the published \code{ibex} configuration; ``one'' is \code{ibex-st}, the same
build with \code{IBEX_CORES=1}. Query identifiers are retained so the exact
expressions can be found in the methodology page and harness.

\begin{table}[htbp]
\centering
\small
\begin{tabular}{@{}lrrr@{}}
\toprule
Query & 1M, default & 16M, default & 16M, one\\
\midrule
\code{update_price_x2} & 0.212 & 4.363 & 18.995\\
\code{mean_by_symbol} & 0.729 & 7.424 & 30.339\\
\code{sort_price} & 21.431 & 516.333 & 697.558\\
\code{order_head_topk} & 2.650 & 39.369 & 39.481\\
\code{tf_rolling_sum_1m} & 2.880 & 53.370 & 50.783\\
\code{tf_rolling_median_1m} & 29.884 & 485.317 & 483.917\\
\code{tf_resample_1m_ohlc} & 4.342 & 98.409 & 96.602\\
\bottomrule
\end{tabular}
\caption{Selected means from the identified September 2026 benchmark
snapshot~\cite{ibexBenchmarks}. M denotes one million input rows.}
\label{tab:cost-snapshot}
\end{table}

Three observations are useful for understanding the engine. First, at 16M
rows, \code{mean_by_symbol} emits 252 groups, \code{order_head_topk} emits
100 rows, and \code{sort_price} emits all 16M rows. Their timings concern
different work and result sizes. In particular, top-k and full sorting are
not interchangeable benchmark answers: the harness requests the top 100
descending prices for the former and a full ascending sort for the latter.

Second, comparing the reported 16M means gives about $4.35\times$ acceleration
for doubling prices and $4.09\times$ for mean by symbol, but only about
$1.35\times$ for full price sorting and essentially none for this top-k
query. These ratios describe the complete operations at that scale. They do
not mean that every phase has that speedup, or that eight vCPUs are eight
independent physical cores.

Third, the time-series rows expose substantial differences within the same
workload family: the one-minute rolling median takes 485.317 ms at 16M rows,
versus 53.370 ms for the rolling sum. The default and one-core means are close
for these examples and for resampling. That is evidence of little observed
threading benefit in this snapshot, not a statistical verdict about their
small differences. The harness constructs timestamps at one-second spacing
and the expressions include \code{as_timeframe}; these numbers therefore do
not isolate window-state maintenance. Varying timestamp density or window
duration would change a different dimension from simply adding more rows.
The detailed window algorithms remain a later investigation.

\evidence{The source CSV's \code{rows} and \code{avg_ms} columns;
\src{tools/ibex_bench.cpp}, definitions of \code{order_head_topk},
\code{sort_price}, and the TimeFrame benchmark group. The published page also
notes that its cross-engine rolling-EWMA comparison uses different mathematics;
that row is not treated here as a comparison of equivalent computations.}

\subsection{Wall time, CPU work, and memory answer different questions}
Parallel execution overlaps work, so adding per-operator or per-worker times
does not generally reconstruct elapsed query time. The critical dependencies
include build-before-probe, aggregate finalization, and ordered output delivery.
Task submission, buffering, and merging also cost time. A larger thread budget
is useful only where it reduces elapsed time enough to offset those costs.

The repository's \src{benchmarking/work_multiplication.py} complements wall
time with total child-process CPU consumption at different core counts.
\src{MEASURING.md} records why this matters: a historical sharded decoder
repeated work while still showing attractive parallel occupancy. High CPU
growth is a diagnostic signal, not proof of redundant decoding; contention
can also increase CPU time. Counting rows, calls, or bytes is needed to
distinguish those explanations.

The benchmark's memory column is absolute peak resident memory, including
already-resident input and other live process storage. It is not an
operator's incremental allocation count. The Ibex harness attempts to reset
Linux's high-water mark before timed iterations; if that mechanism is
unavailable, the code documents a lifetime-peak fallback. Consequently, these
figures describe the footprint of the recorded harness execution and cannot
be substituted directly for the per-operator state accounting above.

\evidence{\src{MEASURING.md}, work multiplication and profile interpretation;
\src{benchmarking/work_multiplication.py}; \src{tools/ibex_bench.cpp},
\code{reset_peak_rss} and \code{peak_rss_mb}.}

\subsection{Using the wider corpus}
The scale suite provides broad coverage, while the PDS-H query suite,
operator profiles, and focused time-series experiments provide complementary
evidence. Their measurement boundaries and datasets should be recorded before
combining results. Existing measurements are the starting point; new runs are
needed only for a specific unresolved claim, such as the contribution of a
particular decoder or the cost of maintaining an ordering guarantee.

For causal claims about a change, the repository's paired, interleaved A/B
workflow and per-operator profiles are more appropriate than differences
between independently published snapshots. This document has reused the
archived evidence without running a benchmark. It neither extrapolates those
times to other machines nor treats a single snapshot as a current performance
guarantee.
